
%% -------------------------------------------------------------------
\subsection{산업별 권장 파라미터}
\label{subsec:industry-guidelines}
%% -------------------------------------------------------------------

$\DRTO$ 프레임워크의 파라미터는 산업별 특성에 따라 교정(calibration)이
필요하다. \chref{ch:results}의 시뮬레이션 결과, 전문가 의견, 그리고 연구모형의 산업
시나리오 설계를 종합하여, 6개 주요 산업에 대한 권장 파라미터 설정을
\p{[}표~\vref{tab:ch6-industry-params}\p{]}에 제시한다.

\begin{table}[htbp]
\centering
\caption{6개 산업별 $\DRTO$ 권장 파라미터}
\label{tab:ch6-industry-params}
\small
\begin{tabularx}{\textwidth}{l c c c c X}
\toprule
\textbf{산업} & \textbf{$R_0$ (h)} & \textbf{$R_{\max}$ (h)} & \textbf{$O_{\text{raw}}$ 범위} & \textbf{권장 조건} & \textbf{적용 근거} \\
\midrule
제조업       & 4.0 & 16.0 & 0.50--0.70 & C2
  & 생산 설비 고장은 가우시안 분포 적합, 생산 계획 연동으로 $R_0$ 엄격 관리 필요 \\
\addlinespace
금융         & 1.0 & 4.0  & 0.70--0.85 & C3
  & 규제 준수 요건(RTO $\leq$ 2h) 엄격, 시장 변동$\cdot$사이버 공격의 파레토 특성 반영 필수 \\
\addlinespace
의료         & 2.0 & 8.0  & 0.60--0.75 & C3
  & 환자 안전이 최우선, 의료기기 장애, 감염 확산 시나리오의 파레토 위험 대비 \\
\addlinespace
IT 서비스    & 0.5 & 2.0  & 0.75--0.90 & C2
  & SLA 기반 복구, 인프라 장애의 대다수가 가우시안 근사 가능, 높은 관측성 인프라 \\
\addlinespace
물류         & 6.0 & 24.0 & 0.45--0.60 & C2
  & 공급망 연속성, 계절 변동$\cdot$물류 지연은 정규분포 근사, 가장 긴 허용 RTO \\
\addlinespace
소매         & 3.0 & 12.0 & 0.50--0.65 & C2
  & 매출 손실 최소화, 일상적 IT 장애 중심, POS$\cdot$결제 시스템 관측성 활용 \\
\bottomrule
\end{tabularx}
\end{table}

\chg{[표 5-15]의 공통 파라미터는 다음과 같다. 곡률 매개변수는 $\kappa = 1.2$로 \jrev{\citet{lee2026drto}}에서 확정된
매개변수이며, 본 연구는 \frev{분포 환경 적합성} 종합 명제로 \frev{운영 적합성을 시사한다}(\secref{sec:ch4_hypothesis_summary}
참조). 관측성 가중치는 권장 범위 $k_O \in [0.4, 0.6]$이고, C2는 가우시안 불확실성 분포, C3는 파레토
불확실성 분포이며, $R_0 / R_{\max}$은 연구모형 기본 설정으로 1:4 비율을 유지한다. $O_{\text{raw}}$ 범위는 해당 산업의
전형적 모니터링 성숙도를 반영한다.}

\p{[}그림~\vref{fig:ch6-industry-summary}\p{]}는 6개 산업의 핵심 파라미터와
권장 조건을 종합적으로 시각화한 것이다.

\begin{figure}[htbp]
\centering
\resizebox{\textwidth}{!}{%
\begin{tikzpicture}[
    indbox/.style={draw, rounded corners=5pt, minimum width=30mm, minimum height=22mm,
                   align=center, font=\small, line width=0.6pt},
    modelC2/.style={draw, rounded corners=3pt, fill=blue!10, minimum width=12mm,
                    minimum height=6mm, font=\scriptsize\bfseries},
    modelC3/.style={draw, rounded corners=3pt, fill=red!10, minimum width=12mm,
                    minimum height=6mm, font=\scriptsize\bfseries},
    conn/.style={-{Stealth[length=2mm]}, thick, black},
]

% 산업 박스
\node[indbox, fill=green!8]  (mfg) at (0, 0)
  {\textbf{제조업}\\$R_0 = 4.0$h\\$O = 0.50$--$0.70$};
\node[indbox, fill=blue!8]   (fin) at (3.2, 0)
  {\textbf{금융}\\$R_0 = 1.0$h\\$O = 0.70$--$0.85$};
\node[indbox, fill=red!8]    (hc)  at (6.4, 0)
  {\textbf{의료}\\$R_0 = 2.0$h\\$O = 0.60$--$0.75$};
\node[indbox, fill=purple!8] (it)  at (9.6, 0)
  {\textbf{IT 서비스}\\$R_0 = 0.5$h\\$O = 0.75$--$0.90$};
\node[indbox, fill=orange!8] (log) at (12.8, 0)
  {\textbf{물류}\\$R_0 = 6.0$h\\$O = 0.45$--$0.60$};
\node[indbox, fill=yellow!10](ret) at (16.0, 0)
  {\textbf{소매}\\$R_0 = 3.0$h\\$O = 0.50$--$0.65$};

% 모델 레이블
\node[modelC2] at (0, -1.7) {C2};
\node[modelC3] at (3.2, -1.7) {C3};
\node[modelC3] at (6.4, -1.7) {C3};
\node[modelC2] at (9.6, -1.7) {C2};
\node[modelC2] at (12.8, -1.7) {C2};
\node[modelC2] at (16.0, -1.7) {C2};

% 연결
\draw[dashed, black] (mfg.south) -- (0, -1.4);
\draw[dashed, black] (fin.south) -- (3.2, -1.4);
\draw[dashed, black] (hc.south)  -- (6.4, -1.4);
\draw[dashed, black] (it.south)  -- (9.6, -1.4);
\draw[dashed, black] (log.south) -- (12.8, -1.4);
\draw[dashed, black] (ret.south) -- (16.0, -1.4);

% R₀ 크기 축 (상단)
\draw[thick, -{Stealth[length=2.5mm]}, black] (-1.2, 2.0) -- (17.2, 2.0)
    node[right, font=\scriptsize, black] {$R_0$ 증가};
\node[font=\scriptsize, black, above] at (3.2, 2.0) {짧은 RTO};
\node[font=\scriptsize, black, above] at (12.8, 2.0) {긴 RTO};

% 범례
\node[modelC2] at (5.5, -2.8) {C2};
\node[font=\scriptsize, anchor=west] at (6.2, -2.8) {가우시안 (일상적 장애)};
\node[modelC3] at (10.5, -2.8) {C3};
\node[font=\scriptsize, anchor=west] at (11.2, -2.8) {파레토 (꼬리 위험 필수)};

\end{tikzpicture}
}% end resizebox
\caption{6개 산업 $\DRTO$ 적용 요약 다이어그램}
\label{fig:ch6-industry-summary}
\end{figure}

산업별 조건 선택의 핵심 기준은 \textbf{파레토 위험의 존재 여부}이다.
금융과 의료는 사이버 공격, 시장 급변, 감염 확산 등 극단적 사건이
복구 시간을 급격히 연장할 수 있으므로 C3(파레토)를 권장한다.
C3의 CVaR$_{99}$가 C2 대비 $+24.0\%$ 높은 점은 이러한 산업에서
꼬리 위험을 사전에 반영해야 하는 정량적 근거가 된다.
반면, 제조업, IT 서비스, 물류, 소매는 장애의 대다수가 가우시안
분포로 근사 가능하므로 C2가 효율적이며, 필요 시 극단 시나리오에 한하여
C3를 보완적으로 적용할 수 있다.

분포 유형의 판단 오류는 복구 자원의 배정에서 두 가지 상반된 위험으로 이어진다.
가우시안 환경(C2)에 파레토 수준(C3)의 자원을 배정하는 과다예비(over-provisioning)는
불확실성을 과대 추정한 결과로, 유휴 대기 인력이나 미사용 예비 시스템, 과도한 예비
부품과 같은 자원이 상시 유지되어 비용 효율이 저하되나 서비스 수준 협약(SLA) 위반
위험은 낮다. 예를 들어 정상적 IT 장애 복구에 대해 랜섬웨어 대응 수준의 인력과 예비
시스템을 상시 유지하면 유휴 자원 비용이 불필요하게 증가한다. 반대로 파레토
환경(C3)에 가우시안 수준(C2)의 자원만 배정하는 과소예비(under-provisioning)는
불확실성을 과소 추정한 결과로, 비용은 절감되나 극단적 사건 발생 시 SLA 위반과
업무 중단 장기화 위험이 급증한다. 공급망 복합 장애 가능성이 있는 물류 기업이 단순
장비 교체 수준의 예비만 확보하면, P99 수준의 극단적 지연(약 $13.2$시간 이상)이
발생할 때 복구 역량이 부족하여 업무 연속성이 심각하게 위협받는다. 따라서 산업별
조건 선택은 단순한 분포 가정의 문제가 아니라 자원 배정의 비용과 위험을 균형 짓는
의사결정이며, 파레토 위험이 존재하는 산업에서 C3를 권장하는 근거가 여기에 있다.

각 산업에 대한 세부 적용 지침은 다음과 같다.

\chg{제조업의 경우,}
생산 설비 고장에 의한 업무 중단이 주요 위험이며,
관측성은 PLC(Programmable Logic Controller), SCADA, MES(Manufacturing
Execution System) 등의 기존 인프라를 활용한다.
$R_0 = 4.0$h는 일반적인 교대제 생산 라인의 복구 목표를 반영하며,
$O = 0.50$--$0.70$은 중간 수준의 모니터링 성숙도에 해당한다.

\chg{금융의 경우,}
금융감독 당국의 IT 복구 시간 규제(예: RTO $\leq$ 2시간)가 엄격하며,
\chg{사이버 공격, 시장 급변, 결제 시스템 장애}의 파레토 특성을 고려하여
C3를 권장한다. $O = 0.70$--$0.85$의 높은 관측성은 실시간 거래
모니터링 시스템(FDS 등)을 반영한다.

\chg{의료의 경우,}
환자 안전이 최우선이며, 의료기기 장애, 전자 의무기록(EMR) 시스템 중단,
감염병 확산 시나리오에서 복구 지연이 생명에 직결된다.
$R_0 = 2.0$h는 응급 의료 시스템의 전형적 복구 목표이며,
C3 적용으로 극단 시나리오에 대한 사전 대비를 강화한다.

\chg{IT 서비스의 경우,}
SLA(Service Level Agreement) 기반의 엄격한 복구 목표를 가지며,
$R_0 = 0.5$h는 클라우드/SaaS 서비스의 전형적 RTO이다.
$O = 0.75$--$0.90$의 높은 관측성은 APM, 로그 분석, 분산 추적(distributed
tracing) 등 성숙한 모니터링 인프라를 반영한다.

\chg{물류의 경우,}
공급망의 연속성이 핵심이며, 계절 변동, 운송 지연, 창고 시스템 장애가
주요 위험이다. $R_0 = 6.0$h는 본 연구의 기본 설정과 동일하며,
$O = 0.45$--$0.60$은 물류 산업의 상대적으로 낮은 IT 관측성 성숙도를 반영한다.

\chg{소매의 경우,}
POS(Point of Sale) 시스템, 재고 관리, 전자상거래 플랫폼의 가용성이
매출에 직접 영향을 미친다. $R_0 = 3.0$h, $O = 0.50$--$0.65$는
중간 규모 소매 환경의 전형적 특성을 반영하며, C2 기반의 효율적 복구를 권장한다.


\chg{이상의 산업별 권장 파라미터를 회귀식에 대입한 구체적 $\DRTO$ 산출 예시는 회귀식 기반 산출 도구를
다루는 \subsecref{subsec:practical-regression-tool}에 제시한다.}


